\chapter{The platform abstraction}
\label{ch:platform}

\begin{aside}
A TUI library that supports POSIX, macOS, and Windows lives or
dies by its platform abstraction. The OS differences are real:
POSIX uses \code{termios} and \code{epoll}; macOS uses
\code{kqueue}; Windows uses the Console API and
\code{WaitForMultipleObjects}. \maya{} hides all of this behind
a small set of platform classes that satisfy the same concept.
This chapter walks the shape of that abstraction.
\end{aside}

\section{The shape of the problem}

To run a TUI on three platforms, \maya{} must do five things
portably:

\begin{enumerate}[leftmargin=*]
  \item Open the terminal and put it in raw mode.
  \item Read keystrokes and mouse events as they arrive.
  \item Receive signals (window resize, interrupt, shutdown).
  \item Wake the main loop from a worker thread.
  \item Restore the terminal when the program exits.
\end{enumerate}

POSIX, macOS, and Windows each spell these five things
differently. \maya{} pins down the spelling with a concept and
provides three implementations.

\begin{funfact}
The first portable terminal library was \code{termcap} (1978),
which abstracted differences between physical terminal models
(VT100 vs ADM-3a). Modern terminal emulators ironed out most of
that diversity. The remaining portability story is between
\emph{operating systems}, not terminals --- a different axis.
\end{funfact}

\section{The Terminal concept}

\file{maya/include/maya/platform/concepts.hpp} declares:

\begin{lstlisting}
template <class T>
concept Terminal = requires (T& t, std::span<std::byte> buf) {
    { t.enter_raw_mode() }   -> std::same_as<Expected<void, IoError>>;
    { t.leave_raw_mode() }   -> std::same_as<void>;
    { t.read(buf) }          -> std::same_as<Expected<size_t, IoError>>;
    { t.write(buf) }         -> std::same_as<Expected<size_t, IoError>>;
    { t.size() }             -> std::same_as<Expected<TermSize, IoError>>;
    { t.caps() }             -> std::same_as<const TerminalCaps&>;
};
\end{lstlisting}

Six methods. Two for setup/teardown, two for I/O, one for size
queries, one for capabilities. Each platform implementation
satisfies the concept.

\subsection*{TermSize and TerminalCaps}

\begin{lstlisting}
struct TermSize {
    int rows, cols;
    int pixel_width = 0, pixel_height = 0;  // optional
};

struct TerminalCaps {
    bool truecolor = false;
    bool osc8_hyperlinks = false;
    bool kitty_keyboard = false;
    bool sync_update = false;
    int max_colors = 16;
};
\end{lstlisting}

\code{TermSize} comes from \code{ioctl(TIOCGWINSZ)} on POSIX,
\code{GetConsoleScreenBufferInfo} on Win32. \code{TerminalCaps}
is detected at startup by sniffing environment variables and
sometimes issuing terminal queries.

\section{The POSIX backend}

\file{maya/include/maya/platform/posix/} defines the POSIX
\code{Terminal}. The core is wrapping \code{termios}.

\subsection*{Raw mode, in detail}

\begin{lstlisting}
Expected<void, IoError> PosixTerminal::enter_raw_mode() {
    if (tcgetattr(fd_, &saved_) != 0)
        return std::unexpect, IoError{from_errno(errno), "tcgetattr"};

    termios raw = saved_;
    raw.c_iflag &= ~(IGNBRK | BRKINT | PARMRK | ISTRIP
                   | INLCR  | IGNCR  | ICRNL  | IXON);
    raw.c_oflag &= ~OPOST;
    raw.c_lflag &= ~(ECHO | ECHONL | ICANON | ISIG | IEXTEN);
    raw.c_cflag &= ~(CSIZE | PARENB);
    raw.c_cflag |= CS8;
    raw.c_cc[VMIN]  = 1;
    raw.c_cc[VTIME] = 0;

    if (tcsetattr(fd_, TCSAFLUSH, &raw) != 0)
        return std::unexpect, IoError{from_errno(errno), "tcsetattr"};

    raw_active_ = true;
    return {};
}
\end{lstlisting}

Each flag mask disables one piece of cooked behaviour. The
incantations come from W. Richard Stevens' \emph{Advanced
Programming in the UNIX Environment}; \maya{} just copies the
recipe.

\subsection*{Event source: epoll}

The POSIX backend uses \code{epoll} on Linux:

\begin{lstlisting}
PosixEventSource::PosixEventSource() {
    epfd_ = epoll_create1(EPOLL_CLOEXEC);
    epoll_ctl(epfd_, EPOLL_CTL_ADD, stdin_fd_,
              &epoll_event{.events = EPOLLIN, .data = {.fd = stdin_fd_}});
    epoll_ctl(epfd_, EPOLL_CTL_ADD, signal_fd_,
              &epoll_event{.events = EPOLLIN, .data = {.fd = signal_fd_}});
    epoll_ctl(epfd_, EPOLL_CTL_ADD, wake_fd_,
              &epoll_event{.events = EPOLLIN, .data = {.fd = wake_fd_}});
}

Event PosixEventSource::wait(Duration timeout) {
    epoll_event ev;
    int n = epoll_wait(epfd_, &ev, 1, to_ms(timeout));
    if (n == 0) return Event::Timeout{};
    if (ev.data.fd == stdin_fd_)   return drain_stdin();
    if (ev.data.fd == signal_fd_)  return drain_signal_fd();
    if (ev.data.fd == wake_fd_)    return Event::Wakeup{};
    return Event::Unknown{};
}
\end{lstlisting}

Three descriptors wait at once: stdin, signalfd, and a self-pipe
for cross-thread wakeups. Whichever fires first wins.

\subsection*{Signal handling}

POSIX signals are delivered out-of-band. \maya{} uses
\code{signalfd} on Linux to turn signals into file-descriptor
events:

\begin{lstlisting}
sigset_t mask;
sigemptyset(&mask);
sigaddset(&mask, SIGINT);
sigaddset(&mask, SIGTERM);
sigaddset(&mask, SIGWINCH);
sigprocmask(SIG_BLOCK, &mask, nullptr);
int sfd = signalfd(-1, &mask, SFD_NONBLOCK | SFD_CLOEXEC);
\end{lstlisting}

The signals are blocked from default delivery and routed
through \code{sfd}. When \code{epoll\_wait} reports activity on
\code{sfd}, a \code{read} drains the pending signals.

On BSD/macOS, \code{kqueue} has native \code{EVFILT\_SIGNAL}
support that does the same thing without the file-descriptor
indirection.

\subsection*{Wake fd}

\code{wake\_fd\_} is a self-pipe used by background threads to
wake the main loop. A worker thread that finished a task writes
one byte to \code{wake\_fd\_}; the main loop's epoll fires; the
loop drains the pipe and dispatches the queued message.

\begin{aside}
The self-pipe trick is one of the oldest patterns in Unix
network programming. It works because reading from an fd is
\emph{the} primitive the OS gives you for ``wake up when
something happens''.
\end{aside}

\section{macOS backend}

\file{maya/include/maya/platform/macos/} reuses much of the
POSIX backend (the \code{termios} flags are the same) but
overrides the event source.

macOS uses \code{kqueue} natively. The event source registers
stdin, signal filters, and a Mach port for cross-thread wakeups.

There is a quirk: macOS's \code{kqueue} delivers
\code{EVFILT\_SIGNAL} events with edge semantics, so the count
of pending signals must be tracked separately. \maya{}
maintains a small atomic counter for each watched signal.

\section{Win32 backend}

Windows is the outlier.

\subsection*{The Console API}

There is no \code{termios}. Raw mode is set via
\code{SetConsoleMode}:

\begin{lstlisting}
DWORD mode;
GetConsoleMode(hStdin, &mode);
mode |= ENABLE_VIRTUAL_TERMINAL_INPUT;
mode &= ~(ENABLE_PROCESSED_INPUT | ENABLE_LINE_INPUT | ENABLE_ECHO_INPUT);
SetConsoleMode(hStdin, mode);

GetConsoleMode(hStdout, &mode);
mode |= ENABLE_VIRTUAL_TERMINAL_PROCESSING;
mode |= DISABLE_NEWLINE_AUTO_RETURN;
SetConsoleMode(hStdout, mode);
\end{lstlisting}

The \code{ENABLE\_VIRTUAL\_TERMINAL\_INPUT} flag is what makes
Windows 10+ consoles emit VT escape sequences for keystrokes
instead of \code{INPUT\_RECORD}s. \maya{} relies on this; older
Windows versions are not supported.

\subsection*{Event source: WaitForMultipleObjects}

Win32's I/O completion ports are very different from epoll. The
\maya{} backend uses a hybrid: \code{WaitForMultipleObjects} on
the console input handle plus a manual-reset event for
cross-thread wakeup:

\begin{lstlisting}
HANDLE handles[] = { hStdin_, hWake_ };
DWORD r = WaitForMultipleObjects(2, handles, FALSE, timeout_ms);
if (r == WAIT_OBJECT_0) return read_stdin();
if (r == WAIT_OBJECT_0 + 1) {
    ResetEvent(hWake_);
    return Event::Wakeup{};
}
\end{lstlisting}

Simpler than epoll, less scalable, but the console input is the
only ``socket'' we need.

\subsection*{Signal handling}

Windows has fewer signals than POSIX. The main ones \maya{}
cares about:

\begin{itemize}[leftmargin=*]
  \item \code{CTRL\_C\_EVENT} and \code{CTRL\_BREAK\_EVENT},
    delivered via a console control handler.
  \item Window resize events, delivered as
    \code{WINDOW\_BUFFER\_SIZE\_EVENT} on the console input
    handle.
\end{itemize}

The control handler runs on a separate thread, so it writes to
the wake event and the main loop picks up the message.

\section{The select header}

\file{maya/include/maya/platform/select.hpp} picks the right
backend at compile time:

\begin{lstlisting}
#if defined(_WIN32)
using Terminal      = Win32Terminal;
using EventSource   = Win32EventSource;
#elif defined(__APPLE__)
using Terminal      = MacTerminal;
using EventSource   = MacEventSource;
#else
using Terminal      = PosixTerminal;
using EventSource   = PosixEventSource;
#endif
\end{lstlisting}

The rest of \maya{} writes \code{maya::platform::Terminal} and
gets the right type.

\subsection*{Why not virtual functions?}

We considered defining \code{ITerminal} with virtual methods and
selecting at runtime. We rejected it because:

\begin{itemize}[leftmargin=*]
  \item The choice is fixed at compile time anyway --- you do
    not target multiple platforms in one binary.
  \item Virtual dispatch on every read/write adds overhead.
  \item Concepts give better error messages when a method is
    missing.
\end{itemize}

\section{Cross-platform pitfalls}

\begin{warning}
\textbf{Line endings.} POSIX terminals send LF (U+000A) for
Enter; Windows sends CRLF (U+000D U+000A). \maya{} normalises in
the input parser, but if you bypass the parser you will see the
raw bytes.
\end{warning}

\begin{warning}
\textbf{Path separators in errors.} POSIX uses forward slashes,
Win32 uses backslashes. The \code{IoError::context} string
preserves whatever the OS gave us. If you display paths to users,
normalise.
\end{warning}

\begin{warning}
\textbf{Console code pages.} Older Windows consoles default to
code page 437 or 850, not UTF-8. \maya{} calls
\code{SetConsoleOutputCP(CP\_UTF8)} on startup. If you bypass
this, your Unicode output will be mojibake.
\end{warning}

\begin{warning}
\textbf{ConPTY vs raw console.} Windows Terminal hosts a
ConPTY (pseudo-terminal) layer that translates between VT
escapes and the legacy console. Most behaviour is identical to
Linux/macOS, but a handful of escapes (cursor shape, mouse
modes) behave subtly differently. If you see a glitch only on
Windows Terminal, check the escape sequence semantics.
\end{warning}

\section*{Workshop}

\begin{enumerate}[leftmargin=*]
  \item Open \file{maya/include/maya/platform/posix/terminal.hpp}
    and find the \code{enter\_raw\_mode} function. Match each
    flag mask to a line in \emph{stty -a} output.
  \item Read \file{maya/include/maya/platform/select.hpp} and
    confirm which backend is selected on your machine.
  \item On Linux, run \code{strace -e epoll\_wait,read,write}
    on a \maya{} app for one second. Note that almost all time
    is spent blocked in \code{epoll\_wait}.
\end{enumerate}

\begin{recap}
\maya{}'s platform layer is a concept (\code{Terminal}) and
three implementations. The POSIX backend uses \code{termios} and
\code{epoll}; macOS uses \code{kqueue}; Win32 uses
\code{SetConsoleMode} and \code{WaitForMultipleObjects}. The
header \code{platform/select.hpp} picks the right backend by
\code{\#ifdef}, so the rest of \maya{} sees one type.
\end{recap}
